\section{Hash Collisions}\label{sec:4}

MD5 was widely deployed when the original paper was written. The collision
attacks reported in 2004--2005 \cite{ref30,ref31} demonstrate why deployment
and statistical behavior do not establish cryptographic collision resistance. RFC 6151 explicitly rejects MD5 for
uses that require collision resistance \cite{ref35}.
The original statement that no practical full SHA-1 collisions were known
was historical: a practical full SHA-1 collision was published in 2017
\cite{ref34}. SHA-256 has a different design and output size; the absence
of an attack in a historical survey is not a proof of security.

\subsection{Practical Effects of Induced Collisions}\label{sec:4.1}
Accidental collisions and deliberately constructed collisions require
separate models. A small birthday probability for ordinary content says
nothing about an adversary's ability to choose colliding inputs. The
calculations below apply to distinct, non-adversarial contents whose hashes
are modeled as independent uniform values. Repeated identical content is
not a collision between distinct messages.

A collision in an address also does not, by itself, establish that an
application can overwrite data: authorization, duplicate detection, and
write-path semantics must be assessed separately.

\subsection{Non-induced Collisions}\label{sec:4.2}
For an ideal $n$-bit output, $q$ independent trials against one fixed digest
succeed with probability $1-(1-2^{-n})^q$. The mean waiting time is $2^n$
trials, and the 50\% point is approximately $(\ln2)2^n$. For a fixed
original message, finding a different message with that digest is a
\emph{second-preimage} problem. The original $2^{127}$ average for a
128-bit digest was incorrect under this model.

If any of $r$ distinct target digests is acceptable, a fresh independent
trial succeeds with probability $r/2^n$, so its mean waiting time is
$2^n/r$. This generic multi-target model neither gives control over which
file is matched nor proves an attack cost for real MD5 or a compound CA.
Candidate trials can be evaluated offline; they need not be stored in a
cluster. Any-pair collisions instead follow the birthday calculation in
\S\ref{sec:5.1}.
